\chapter{System Architecture}
\label{ch:architecture}

\section{Layering}

Xyra is organised as five layers. Each layer may call downward and emit events upward. No
layer may call upward directly. This rule is what keeps the surfaces replaceable and the
core testable without a display.

\begin{figure}[h]
\centering
\begin{tikzpicture}[
  layer/.style={draw=xyraline,fill=xyrapanel,minimum width=118mm,minimum height=13mm,rounded corners=2pt,font=\small},
  lab/.style={font=\scriptsize\color{xyragrey}}
]
\node[layer] (l5) at (0,0) {\textbf{Surfaces} \quad Workbench, Mission Control, Agent Theatre, Inspector, Deck, Atlas};
\node[layer,below=3mm of l5] (l4) {\textbf{Orchestration} \quad Mission supervisor, council, scheduler, budget, journal};
\node[layer,below=3mm of l4] (l3) {\textbf{Intelligence} \quad Context engine, compression, retrieval, graph, skills};
\node[layer,below=3mm of l3] (l2) {\textbf{Workspace} \quad Buffers, syntax, language servers, git, terminal, filesystem};
\node[layer,below=3mm of l2] (l1) {\textbf{Platform} \quad UI framework, async runtime, storage, process supervision, sandbox};

\draw[-{Stealth[length=2mm]},xyralime,line width=0.8pt] ($(l1.west)+(-6mm,0)$) -- ($(l5.west)+(-6mm,0)$) node[midway,left,lab,rotate=90,anchor=south] {events upward};
\draw[-{Stealth[length=2mm]},xyragrey,line width=0.8pt] ($(l5.east)+(6mm,0)$) -- ($(l1.east)+(6mm,0)$) node[midway,right,lab,rotate=-90,anchor=south] {calls downward};
\end{tikzpicture}
\caption{Five layer architecture with a strict dependency direction.}
\end{figure}

\section{Process model}

Xyra runs as a small constellation of processes rather than one monolith. The reason is
containment: an agent process that hangs, a language server that leaks, or an extension that
panics must never take down the editor, and must never be able to block the interface
thread.

\begin{figure}[h]
\centering
\begin{tikzpicture}[
  proc/.style={draw=xyraline,fill=xyrapanel,rounded corners=2pt,minimum height=10mm,minimum width=34mm,font=\small},
  main/.style={draw=xyralime,line width=1pt,fill=white,rounded corners=2pt,minimum height=12mm,minimum width=40mm,font=\small\bfseries},
  pipe/.style={-{Stealth[length=2mm]},xyragrey}
]
\node[main] (ui) at (0,0) {xyra (UI process)};
\node[proc,below=14mm of ui] (core) {xyrad (core service)};
\node[proc,left=16mm of core] (ls) {language servers};
\node[proc,right=16mm of core] (ag) {agent processes};
\node[proc,below left=12mm and 4mm of core] (sb) {sandbox runners};
\node[proc,below right=12mm and 4mm of core] (ext) {WASM extensions};

\draw[pipe] (ui) -- node[right,font=\scriptsize] {local socket, typed messages} (core);
\draw[pipe] (core) -- node[above,font=\scriptsize] {LSP} (ls);
\draw[pipe] (core) -- node[above,font=\scriptsize] {ACP} (ag);
\draw[pipe] (core) -- node[left,font=\scriptsize,xshift=-1mm] {process} (sb);
\draw[pipe] (core) -- node[right,font=\scriptsize] {host calls} (ext);
\end{tikzpicture}
\caption{Process topology. The interface never speaks to an agent or a language server directly.}
\end{figure}

\begin{description}[style=nextline,leftmargin=0pt]
  \item[xyra] The interface process. Owns the window, the element tree, input handling and
  rendering. Holds no authoritative state; it renders a projection pushed from the core and
  sends intents back.

  \item[xyrad] The core service. Owns buffers, the project model, the context engine, the
  mission supervisor, the journal and all persistence. Survives interface restarts, which
  means a crashed window never loses a running mission.

  \item[Language servers] One child per language per workspace, supervised with restart
  backoff and a health probe.

  \item[Agent processes] External agent binaries speaking the agent client protocol over
  stdio. Each runs in its own working directory with an environment the core constructs.

  \item[Sandbox runners] Ephemeral processes executing tests inside an isolated worktree.
  Killed on budget expiry without affecting anything else.

  \item[WASM extensions] Untrusted third party code with an explicit capability list, no
  ambient filesystem access and no network unless granted.
\end{description}

\section{Crate topology}

\begin{longtable}{@{}p{0.26\textwidth} p{0.66\textwidth}@{}}
\toprule
\textbf{Crate} & \textbf{Responsibility} \\
\midrule
\endfirsthead
\toprule
\textbf{Crate} & \textbf{Responsibility} \\
\midrule
\endhead
xyra-ui & Element model wrapper over the vendored UI framework, theme tokens, primitives \\
xyra-text & Rope buffer, transactions, undo history, multi cursor selection algebra \\
xyra-syntax & Parser host, incremental reparse, highlight queries, structural selection \\
xyra-lang & Language server client, capability negotiation, diagnostics, completion \\
xyra-project & Worktree model, ignore semantics, file watching, path resolution \\
xyra-vcs & Git operations, status, diff, blame, staging, commit, branch \\
xyra-term & Terminal grid, PTY supervision, escape parsing, search \\
xyra-context & Symbol and import graph, retrieval, compression pipeline, budget solver \\
xyra-memory & Durable stores, caches, journal, mission state, migrations \\
xyra-agent & Agent process supervision, protocol clients, tool exposure, transcript capture \\
xyra-mission & Ticket graph, supervisor loop, failure policy, budgets \\
xyra-verify & Sandbox execution, visual verification, quality assurance sweeps, flags \\
xyra-fleet & Multi repository registry, cross repository search and impact \\
xyra-skills & Skill registry, resolution, injection, precedence \\
xyra-ext & Extension host, capability enforcement, manifest validation \\
xyra-proto & Message types shared between the interface and the core \\
xyrad & Core service binary, wiring, supervision tree \\
xyra & Interface binary, surfaces, keymap, command registry \\
\bottomrule
\caption{Crate responsibilities. One crate, one reason to change.}
\end{longtable}

\section{State ownership and the projection model}

The interface holds no authoritative state. Every surface receives a projection: a
serialisable snapshot of exactly what it needs to render, versioned so that stale frames are
discarded. User actions travel back as intents, never as mutations.

\begin{figure}[h]
\centering
\begin{tikzpicture}[
  box/.style={draw=xyraline,fill=xyrapanel,rounded corners=2pt,minimum height=9mm,minimum width=30mm,font=\small},
  arr/.style={-{Stealth[length=2mm]},xyragrey}
]
\node[box] (intent) {Intent};
\node[box,right=16mm of intent] (core) {Core reducer};
\node[box,right=16mm of core] (state) {Authoritative state};
\node[box,below=11mm of state] (proj) {Projection};
\node[box,left=16mm of proj] (surface) {Surface render};

\draw[arr] (intent) -- (core);
\draw[arr] (core) -- (state);
\draw[arr] (state) -- (proj);
\draw[arr] (proj) -- (surface);
\draw[arr] (surface) to[out=180,in=180,looseness=1.4] node[left,font=\scriptsize] {user action} (intent);
\end{tikzpicture}
\caption{Unidirectional flow. The loop is closed only by user action.}
\end{figure}

This is what makes the product testable. Every surface can be exercised by feeding
projections and asserting on the intents produced, with no window open and no GPU present.

\section{Concurrency model}

\begin{itemize}
  \item The interface thread never blocks. Any operation with unbounded duration is a task
  on the core's runtime, and the interface shows an intermediate state.
  \item Buffer edits are applied on a single owning task per buffer. Concurrent edits arrive
  as transactions and are ordered, never merged optimistically.
  \item Indexing, embedding and graph construction run on a bounded worker pool sized to
  physical cores minus two, so the machine stays usable.
  \item Every long running task carries a cancellation token and a budget. A task without a
  budget is a defect.
  \item Agent sessions are supervised with a timeout, an inactivity watchdog and a hard kill
  path, because the dominant real world failure is a subprocess that never exits.
\end{itemize}

\section{Persistence}

\begin{table}[h]
\centering
\small
\begin{tabularx}{\textwidth}{l l X}
\toprule
\textbf{Store} & \textbf{Technology} & \textbf{Contents} \\
\midrule
Project index & SQLite & Symbols, import edges, file hashes, chunk metadata \\
Vector index & Embedded ANN file & Chunk embeddings, quantised, memory mapped \\
Mission state & JSON with atomic replace & Ticket graph, attempts, budgets, commits \\
Journal & Append only JSON lines & Decisions, verdicts, verification results \\
Caches & Content addressed files & Compression artefacts, screenshots, transcripts \\
Settings & TOML & User and workspace configuration, layered \\
\bottomrule
\end{tabularx}
\caption{Storage responsibilities. Nothing important lives only in memory.}
\end{table}

Every store is per workspace under \ctl{.xyra/} except global caches, which live under the
platform cache directory. Deleting \ctl{.xyra/} must degrade the product to a cold start and
never corrupt the repository.

\section{Failure containment}

\begin{longtable}{@{}p{0.30\textwidth} p{0.62\textwidth}@{}}
\toprule
\textbf{Failure} & \textbf{Containment} \\
\midrule
\endhead
Language server crash & Supervised restart with exponential backoff, diagnostics preserved, editing unaffected \\
Agent process hang & Inactivity watchdog kills the session, ticket is marked failed, supervisor continues \\
Extension panic & WASM trap is caught, extension is disabled for the session, a flag is raised \\
Core service crash & Interface shows a reconnect state, mission resumes from durable state on restart \\
Interface crash & Core keeps working, mission continues, window reopens onto live state \\
Disk full & Writes fail closed, mission halts with a stated reason, nothing is half written \\
Corrupt index & Detected by checksum, index rebuilt in background, structural queries degrade to search \\
\bottomrule
\caption{Every failure has a named containment strategy.}
\end{longtable}
